\section{Conclusion}
\label{sec:conclusion}

We presented AnomaClaw, a training-free system for cross-domain visual anomaly detection that augments VLM reasoning with structured patch-level expert evidence. Our central finding is that expert evidence is most effective when provided to VLMs as structured textual context---rather than as a routing signal or score-fusion prior---enabling the VLM to ground its reasoning in quantitative patch-level statistics while retaining full decision authority. This expert-as-context design addresses the fundamental lack of quantitative grounding that limits VLM-only approaches, outperforming expert-as-router by 3.3 AUROC points and score-fusion by 3.0 points in matched ablations.

Evaluated on 1{,}200 test items across 10 diverse domains, AnomaClaw achieves 0.882 macro AUROC with family-adaptive fusion, surpassing a pure VLM baseline (0.866) and a multi-round agent baseline (0.873). The per-domain analysis reveals that the expert signal contributes most on local-appearance domains, where patch distances are strongly discriminative, and least on medical domains, where the VLM must override expert evidence due to feature overlap between normal and pathological tissue. The causal gain from structured expert context over generic text context (+0.7 AUROC points) is positive but not statistically significant at the domain level ($p = 0.21$), suggesting that the primary contribution is the \emph{integration strategy} (context vs.\ routing vs.\ fusion) rather than the specific expert content.

Several limitations warrant discussion. Performance on liver CT data (D5c, 0.719 AUROC) remains substantially below practical clinical requirements, reflecting fundamental limits of both DINOv2 patch features on medical imagery and VLM visual understanding for subtle pathology. The system relies on a proprietary VLM (GPT-5.4), constraining reproducibility. The domain knowledge text and family labels used for calibration are hand-crafted; automatic family detection would be needed to extend to previously unseen domains without human intervention. Finally, our benchmark evaluates 120 items per domain; larger-scale evaluation would strengthen statistical power.

Looking forward, replacing the proprietary VLM backbone with capable open-source alternatives (e.g., InternVL, Qwen2.5-VL) would improve accessibility. Domain-specific preprocessing such as CT windowing could address the medical domain weakness. Beyond classification, extending to anomaly localization and integrating active learning represent natural next steps toward a practical cross-domain anomaly detection system.
